% Part: sets-functions-relations
% Chapter: functions
% Section: partial-functions

\documentclass[../../../include/open-logic-section]{subfiles}


\begin{document}

\olfileid{sfr}{fun}{par}

\olsection{جزوي تابعې}

\begin{explain}
کله کله ګټوره وي چې د تابعې په تعريف کښې نرمي راولو، څو د ټولو
ممکنو ننوتونو دپاره د تابعې د وت تعريف کېدل لازم نۀ وي۔ داسې
نگاشتونه \emph{جزوي تابعې} بلل کېږي۔
\end{explain}

\begin{defn}
\emph{جزوي تابع} $f \colon A \pto B$ داسې نگاشت دے چې د~$A$ له
هر \psOblique{!!{element}} سره د~$B$ تر يوه زيات !!{element} نۀ
تړي۔ که $f$ له $x \in A$ سره د~$B$ يو غړے وتړي، وايو چې $f(x)$
\emph{تعريف شوے} دے، او که داسې نۀ وي، \emph{تعريف نۀ شوے} دے۔
که $f(x)$ تعريف شوے وي، $f(x) \fdefined$ ليکو؛ کنه،
$f(x) \fundefined$ ليکو۔ د جزوي تابعې~$f$ \emph{د تعريف ساحه}
د~$A$ هغه فرعي سټ دے چې تابع پرې تعريف شوې ده؛ يعنې
$\dom{f} = \Setabs{x \in A}{f(x) \fdefined}$۔
\end{defn}

\begin{ex}
هره تابع $f\colon A \to B$ يوه جزوي تابع هم ده۔ هغه جزوي تابعې
چې پر~$A$ هرځاے تعريف شوې وي---يعنې هغه څه چې تر دې دمه مو
يوازې تابع بلل---\emph{هرځاے تعريف شوې} تابعې هم بلل کېږي۔
\end{ex}

\begin{ex}
جزوي تابع $f \colon \Real \pto \Real$ چې په $f(x) = 1/x$ ورکړل
شوې، د $x = 0$ دپاره تعريف شوې نۀ ده، او په نورو ټولو ځايونو
کښې تعريف شوې ده۔
\end{ex}

\begin{prob}
د ورکړل شوې $f\colon A \pto B$ دپاره، جزوي تابع $g\colon B \pto
A$ داسې تعريف کړئ: د هر $y \in B$ دپاره، که يو يوازينے
$x \in A$ وي چې $f(x) = y$ وي، نو $g(y) = x$؛ کنه،
$g(y) \fundefined$۔ وښيئ چې که $f$ يو پر يو وي، نو د ټولو
$x \in \dom{f}$ دپاره $g(f(x)) = x$، او د ټولو $y \in \ran{f}$
دپاره $f(g(y)) = y$ وي۔
\end{prob}

\begin{defn}[د جزوي تابعې ګراف]
فرض کړئ $f\colon A \pto B$ يوه جزوي تابع ده۔ د~$f$ \emph{ګراف}
هغه اړيکه $R_f \subseteq A \times B$ ده چې داسې تعريف شوې ده:
\[
R_f = \Setabs{\tuple{x,y}}{f(x) = y}.
\]
\end{defn}

\begin{prop}
فرض کړئ $R \subseteq A \times B$ دا خاصيت لري چې هرکله $Rxy$
او $Rxy'$ وي، نو $y = y'$ وي۔ بيا $R$ د هغې جزوي تابعې
$f\colon A \pto B$ ګراف ده چې داسې تعريف شوې ده: که داسې $y$
وي چې $Rxy$ وي، نو $f(x) = y$؛ کنه، $f(x) \fundefined$۔ که $R$
\emph{مسلسله} هم وي، يعنې د هر $x \in A$ دپاره داسې $y \in B$
وي چې $Rxy$ وي، نو $f$ هرځاے تعريف شوې ده۔
\end{prop}

\begin{proof}
فرض کړئ داسې $y$ شته چې $Rxy$ وي۔ که يو بل $y' \neq y$ هم واے
چې $Rxy'$ واے، پر $R$ لګېدلے شرط به مات شوے و۔ نو که داسې $y$
شته چې $Rxy$ وي، هغه $y$ يوازينے دے؛ له دې کبله $f$ په سمه
توګه تعريف شوې ده۔ په ښکاره، $R_f = R$، او که~$R$ مسلسله وي،
نو $f$ هرځاے تعريف شوې ده۔
\end{proof}

\end{document}
